% !TeX root = ../navier-stokes-manuscript.tex
\subsection{The temporal Tower}\label{sub:BodyFiniteTimeResponseEscalation}

The comparison clocks give a shrinking schedule.
They do not tell us how one velocity field changes across it.
On any strict interval before \(T_*\), the same solution remains classical, so hold \(x\) fixed and write
\begin{equation}\label{eq:TemporalTowerDefinitionBody}
  V_n:=\partial_t^nu,
  \qquad
  Q_n:=\partial_t^np,
  \qquad
  n\in\Nzero.
\end{equation}
Each \(V_n(\cdot,t)\) is still a spatial field.
The differentiation is Eulerian; it does not follow the moving vortex.

The first response is already in \NavierStokes{}.
For \(V_0=u\),
\[
  V_1
  =\nu\Delta V_0
  -(V_0\cdot\nabla)V_0
  -\nabla Q_0.
\]
Differentiate once.
The derivative can land on either copy of velocity in the transport term:
\[
  V_2
  =\nu\Delta V_1
  -(V_1\cdot\nabla)V_0
  -(V_0\cdot\nabla)V_1
  -\nabla Q_1.
\]
Differentiate again.
The middle placement occurs twice:
\[
  V_3
  =\nu\Delta V_2
  -(V_2\cdot\nabla)V_0
  -2(V_1\cdot\nabla)V_1
  -(V_0\cdot\nabla)V_2
  -\nabla Q_2.
\]
At order \(n\), the same product rule distributes the time derivatives between the transporting velocity and the velocity being transported:
\begin{equation}\label{eq:BodyTemporalVelocityRecurrence}
  V_{n+1}
  +\sum_{a=0}^{n}\binom{n}{a}(V_a\cdot\nabla)V_{n-a}
  +\nabla Q_n
  =\nu\Delta V_n,
  \qquad
  \nabla\cdot V_n=0.
\end{equation}
These successive rows are the exact temporal Tower generated by \NavierStokes{}.

The recurrence still contains \(Q_n\).
Pressure is not an independent row: the divergence constraint recovers it from the velocity.
Taking the divergence of the original equation and using decay at spatial infinity gives
\begin{equation}\label{eq:PressureNormalization}
  -\Delta p=\partial_i\partial_j(u_i u_j),
  \qquad
  p=R_iR_j(u_i u_j).
\end{equation}
Here \(R_j:=\partial_j(-\Delta)^{-1/2}\), and the decaying normalization removes the arbitrary spatial constant in pressure at each time.
After \(n\) time derivatives, the same binomial splitting gives
\begin{equation}\label{eq:BodyTemporalPressureRecurrence}
  Q_n
  =R_iR_j
  \sum_{a=0}^{n}\binom{n}{a}
  (V_a)_i(V_{n-a})_j.
\end{equation}
The velocity rows determine \(Q_n\), and \(\nabla Q_n\) returns in the equation for \(V_{n+1}\).

Temporal order does not identify everything inside a row.
Viscosity contributes \(\Delta V_n\), transport contributes a spatial derivative, and pressure returns through spatial singular integrals of the same velocity products.
Those operations open an independent spatial coordinate:
\begin{equation}\label{eq:BodyMixedJetDefinition}
  J_{n,\alpha}:=\partial_t^n\partial_x^\alpha u,
  \qquad
  \Pi_{n,\alpha}:=\partial_t^n\partial_x^\alpha p,
  \qquad
  \alpha\in\Nzero^3.
\end{equation}
One spatial derivative splits the transport term again:
\[
  \partial_\ell\bigl((V_a\cdot\nabla)V_{n-a}\bigr)
  =((\partial_\ell V_a)\cdot\nabla)V_{n-a}
  +(V_a\cdot\nabla)\partial_\ell V_{n-a}.
\]
For a general multi-index \(\alpha\), Leibniz's rule gives
\begin{align}
  J_{n+1,\alpha}
  &+
  \sum_{a=0}^{n}
  \sum_{\beta\leq\alpha}
  \binom{n}{a}\binom{\alpha}{\beta}
  (J_{a,\beta}\cdot\nabla)
  J_{n-a,\alpha-\beta}
  +\nabla\Pi_{n,\alpha}
  \notag\\
  &=
  \nu\Delta J_{n,\alpha},
  \qquad
  \nabla\cdot J_{n,\alpha}=0.
  \label{eq:BodyMixedJetRecurrence}
\end{align}
The pressure coordinates follow from the same recovery:
\begin{equation}\label{eq:BodyMixedPressureRecurrence}
  -\Delta\Pi_{n,\alpha}
  =
  \partial_x^\alpha\partial_i\partial_j
  \sum_{a=0}^{n}\binom{n}{a}
  (V_a)_i(V_{n-a})_j.
\end{equation}

Each temporal rung has opened an entire spatial family.
These identities show how the coordinates are generated; they do not bound them.
What remains is to count how much spatial height each temporal rung carries.
